%=========================
% Sections/08-KetLuanVaHuongPhatTrien.tex
% Chương 8 -- Kết luận và hướng phát triển
%=========================

\section{Kết luận và hướng phát triển}

Chương này tổng kết kết quả của đề tài và nêu các hướng phát triển trong tương lai.

\subsection{Kết luận}

Đề tài đã xây dựng hệ thống thương mại điện tử thời trang có trợ lý mua sắm đa phương thức. Khách hàng có thể quản lý tài khoản, tìm sản phẩm, sử dụng giỏ hàng, lưu sản phẩm yêu thích, áp dụng mã giảm giá, đặt hàng COD, theo dõi đơn và đánh giá sản phẩm. Quản trị viên có thể quản lý người dùng, sản phẩm, tồn kho, đơn hàng, khuyến mãi và bài viết. Giao diện web sử dụng được trên máy tính và thiết bị di động.

Hệ thống hỗ trợ tìm kiếm bằng văn bản, hình ảnh và kết hợp ảnh--văn bản. Mô hình YOLOv26 đã được huấn luyện thêm trên dữ liệu thời trang để nhận diện và cắt vùng sản phẩm trong ảnh. Trợ lý sử dụng RAG và bộ nhớ hội thoại để tư vấn dựa trên dữ liệu sản phẩm, đồng thời ghi nhớ nội dung của các lượt trao đổi trước.

Hệ thống gồm React Frontend, Spring Boot Backend, FastAPI AI Service, MySQL và Qdrant. MySQL lưu dữ liệu nghiệp vụ chính, còn Qdrant lưu dữ liệu vector cho tìm kiếm. Backend luôn kiểm tra lại giá, tồn kho và trạng thái sản phẩm trước khi tạo giao dịch. Vì vậy, các chức năng mua sắm chính vẫn hoạt động khi AI Service bị gián đoạn. Hệ thống đã được đóng gói bằng Docker, chạy trên máy local và triển khai thử nghiệm trên Google Cloud Platform.

Các chức năng AI đều đạt tiêu chí đã đặt ra. Text Search đạt Hit@5 là 1,000 và NDCG@5 là 0,941, cao hơn BM25. Image Search đạt Exact Hit@5 là 1,000 với ảnh catalog và 0,970 với ảnh bên ngoài. Hybrid Search đạt Strict Hit@5 là 0,980 và tỷ lệ thỏa ràng buộc@5 là 0,990. Trợ lý hội thoại đạt Grounding là 1,000, Context Precision là 0,946, Context Recall là 0,944 và vượt qua 12/12 ca kiểm thử quan trọng. Qua so sánh các cấu hình, hệ thống chọn cấu hình không rerank vì cho kết quả phù hợp nhất.

Phần nghiệp vụ đạt 102/102 ca kiểm thử Backend và hoàn thành năm luồng end-to-end. MySQL và Qdrant đồng bộ đủ 1.000 \texttt{product\_id}; các chức năng thử lại, sao lưu và khôi phục đều hoạt động đúng. Backend xử lý thành công 6.050/6.050 request hỗn hợp và 57.600/57.600 request trong bài kiểm thử kéo dài hai giờ. Các chức năng AI đạt thời gian phản hồi P95 không quá 5 giây trong mức tải đã thử. Các kiểm tra về đăng nhập, phân quyền, dữ liệu đầu vào, giao dịch, xử lý lỗi và theo dõi request cũng đạt yêu cầu.

Nhìn chung, đề tài đã hoàn thành các mục tiêu đề ra. Sản phẩm kết hợp được các chức năng mua sắm chính với tìm kiếm đa phương thức và trợ lý hội thoại. Người dùng có thêm nhiều cách mô tả nhu cầu và tìm sản phẩm phù hợp thay vì chỉ nhập từ khóa.

\subsection{Hướng phát triển}

Trong tương lai, hệ thống có thể được tiếp tục phát triển theo các hướng sau:

\begin{itemize}
    \item \textbf{Mở rộng dữ liệu:} tăng số lượng sản phẩm, bổ sung nhiều ngành hàng, truy vấn và hình ảnh thực tế. Dữ liệu dùng để điều chỉnh hệ thống cần tách riêng với dữ liệu dùng để đánh giá.

    \item \textbf{Cải thiện tìm kiếm:} giúp hệ thống hiểu tốt hơn các yêu cầu về màu sắc, giá, loại sản phẩm và mục đích sử dụng. Phần xử lý ảnh có thể được bổ sung chức năng tách nền, nhận diện vật thể tốt hơn và mô hình chuyên cho sản phẩm thương mại điện tử.

    \item \textbf{Cá nhân hóa gợi ý:} sử dụng lịch sử tìm kiếm, sản phẩm đã xem, danh sách yêu thích và đơn hàng để đưa ra gợi ý phù hợp hơn. Người dùng cần có quyền xem và xóa dữ liệu cá nhân hóa của mình.

    \item \textbf{Bổ sung chức năng mua sắm:} tích hợp thanh toán trực tuyến, đơn vị vận chuyển, theo dõi giao hàng, đổi trả và hoàn tiền. Hệ thống cũng có thể phát triển ứng dụng di động và thông báo theo thời gian thực.

    \item \textbf{Tăng khả năng phục vụ:} kiểm thử với nhiều người dùng và trong thời gian dài hơn. Hệ thống có thể tự tăng tài nguyên khi tải cao, dùng hàng đợi cho tác vụ đồng bộ và bổ sung bản sao dữ liệu để giảm gián đoạn.

    \item \textbf{Tăng bảo mật và theo dõi hệ thống:} thực hiện kiểm thử bảo mật chuyên sâu và diễn tập khôi phục định kỳ. Bảng theo dõi cần hiển thị thời gian phản hồi, tỷ lệ lỗi, tài nguyên sử dụng, chi phí AI và chất lượng tìm kiếm.

    \item \textbf{Đánh giá với người dùng thực:} đo thời gian tìm sản phẩm, tỷ lệ hoàn thành thao tác và mức độ hài lòng. Kết quả sẽ giúp cải thiện giao diện và cách trợ lý đặt câu hỏi khi yêu cầu chưa rõ.
\end{itemize}

Trước mắt, hệ thống nên ưu tiên mở rộng dữ liệu đánh giá và cải thiện tìm kiếm kết hợp. Sau đó cần kiểm thử tải dài hơn, tăng bảo mật và hoàn thiện phần theo dõi trước khi đưa vào sử dụng ở quy mô lớn.
